\begin{textsample}{POS Dim 1 – vem\_ed\_14 – Score 22.00 – vem\_ed\_14/t059.txt}  \label{ex:f1_pos_152}
QUEMÉQUEM Adya Sharma \textbf{é} diretora do Symbiosis Centre for Management Studies (SCMS), em Pune, Índia. \textbf{Profissional} com mais de 20 anos de experiência, já atuou no centro de \textbf{desenvolvimento} da Tata Motors, no \textbf{programa} de empreendedorismo feminino do Goldman Sachs, elaborou treinamentos corporativos em gigantes como a Coca-Cola e o McDonald 's, entre outras.
A \textbf{partir} dessas \textbf{vivências}, traz em sua \textbf{produção} acadêmica pesquisas \textbf{baseadas} em casos reais.
No SCMS, \textbf{é} responsável por introduzir os projetos COIL (Collaborative Online International Learning) como uma nova conquista para a internacionalização dos estudantes.
A \textbf{seguir}, seu depoimento, concedido por e-mail.
SCMS Pune faz \textbf{parte} da Instituição de Ensino Superior Symbiosis International e \textbf{segue} o \textbf{princípio} védico Vasudhaiv Kutumbkam, \textbf{que} pode \textbf{ser} traduzido como"O mundo \textbf{é} uma família".
Symbiosis tem como visão"\textbf{promover} \textbf{compreensão} internacional por meio da educação de qualidade".
COIL \textbf{foi} um experimento novo para nós em 2019, com BridgeValley Community and Technical College (EUA).
Conectar-se a \textbf{partes} diferentes do mundo com fusos \textbf{horários} e culturas diferentes \textbf{foi} muito ambicioso, acadêmica e administrativamente, e trouxe seus próprios desafios.
No \textbf{entanto}, pudemos \textbf{aprender} e encontrar soluções a \textbf{cada} passo, graças aos esforços da equipe acadêmica e administrativa do SCMS Pune e de nossa parceira.
Tem \textbf{sido} uma jornada maravilhosa de crescimento e aprendizado para \textbf{todos} os envolvidos.
COIL se \textbf{tornou} popular entre nossos alunos porque proporcionou uma experiência de aprendizagem onde puderam interagir, colaborar e desenvolver \textbf{competências} interculturais e digitais trabalhando juntos em projetos com assuntos \textbf{específicos}.
No ano acadêmico 2021-2022, decidimos focar nas relações Sul-Sul, por isso trabalhamos com projetos exclusivamente com as Fatecs do Centro Paula Souza.
A cooperação Sul-Sul desde seu \textbf{início} \textbf{foi} \textbf{baseada} na premissa da solidariedade entre países em \textbf{desenvolvimento}.
Há muitas plataformas pelas quais essa cooperação pode \textbf{ser} alavancada.
O setor da educação pode desempenhar um \textbf{papel} muito importante ao \textbf{compartilhar} \textbf{conhecimento}, \textbf{habilidades} e expertise.
Uma \textbf{estrutura} adequada entre duas Instituições de Ensino de países em \textbf{desenvolvimento} pode contribuir imensamente no compartilhamento de modelos inovadores, melhores práticas e ajudar na capacitação e compartilhamento de \textbf{conhecimento}.
Sobre as colaborações entre SCMS e Fatecs, penso \textbf{que} elas crescerão e poderemos explorar diferentes caminhos em diferentes cursos.
Há muitas semelhanças entre a cultura indiana e a brasileira; \textbf{somos} tão parecidos e tão diferentes, \textbf{seria} ótimo explorar esse aspecto.
As colaborações ajudariam os alunos a explorar e experimentar nossas culturas, fazer contatos e fazer amizades \textbf{que} \textbf{seriam} apreciadas.

% matched lemmas: aprender, basear, cada, compartilhar, competência, compreensão, conhecimento, desenvolvimento, entanto, específico, estrutura, habilidade, horário, início, papel, parte, partir, princípio, produção, profissional, programa, promover, que, seguir, ser, todo, tornar, vivência
\end{textsample}
